\section{Introduction}
\label{sec:introduction}

Probabilistic representations for nonlinear parabolic equations begin from a simple algebraic observation. Duhamel's formula converts the nonlinearity into a time integral, and a branching process converts products inside that integral into products over descendants. If the equation contains spatial derivatives of the solution, Gaussian integration by parts or a Malliavin weight transfers those derivatives to the sampled motion. At finite depth this procedure is often exact. The difficulty addressed in this paper is that exactness is a signed statement, whereas a probabilistic representation requires an absolute-moment statement.

This distinction becomes acute for derivative weights. A second derivative of the heat semigroup may be written as
\begin{equation}
  \partial_x^2P_rF(x)
  =
  \E\br{
    \frac{\He_2(Z)}{r}
    \pa{F(x+\sqrt r Z)-F(x)}
  },
  \qquad
  \He_2(z)=z^2-1,
  \label{eq:intro-centered-hessian}
\end{equation}
where $Z$ is standard Gaussian. The centering in~\eqref{eq:intro-centered-hessian} is the source of the regularizing factor. It is also destroyed if the absolute value is taken too early. Thus the main question is not only whether the branching identity is correct, but where cancellation is allowed to occur before the first absolute value.

The results below support the following principle.

\begin{quote}
For branching representations with singular derivative weights, cancellation before absolute values is a structural part of integrability, not a variance-reduction refinement added after the representation has been constructed.
\end{quote}

The paper reaches this conclusion in three stages. The first stage concerns existing branching representations and shows that $L^1$ can depend on the architecture even when the PDE and terminal datum are fixed. The second stage constructs an integrable representation for a quadratic Hessian equation by averaging complete patch interiors before absolute values. The third stage proves that, for one fixed smooth datum, this averaging cannot be replaced by a proposal change which keeps the canonical raw signed marked contribution as its conditional barycenter.

\subsection{Formal exactness and the coding-tree obstruction}

Nguwi, Penent, and Privault construct a coding-tree Feynman--Kac formula for fully nonlinear equations in which the nonlinearity may depend on derivatives of arbitrary order~\cite{NguwiPenentPrivault2023}. Their code set is closed under the differentiations generated by the PDE. In particular, a composite code contains a Hessian mechanism which may be repeated indefinitely along one distinguished lineage.

For the terminal problem
\begin{equation}
  \partial_tu+\frac12u_{xx}+f(J_nu)=0,
  \qquad
  u(T)=\phi,
  \label{eq:intro-general-pde}
\end{equation}
we isolate this repeated Hessian genealogy. If $g^*$ is an allowed composite code, $j$ is an active jet direction, and
\begin{equation*}
  D_m(B;g,j)
  =
  \int_B
  \abs{\partial_{z_j}^{2m}g(J_n\phi(y))}\,dy,
\end{equation*}
then the contribution of the length-$m$ restricted genealogy has the lower-bound scale
\begin{equation*}
  c^m\frac{D_m(B;g,j)}{m!}.
\end{equation*}
The lifetime density, the mechanism-selection probability, and the survival probabilities cancel their reciprocal compensators. The only combinatorial loss is the ordered-time simplex. This gives the repeated-Hessian obstruction, stated as~\cref{thm:repeated-hessian}. In particular, if $D_m/m!$ has superexponential growth, the raw code-rooted functional is not in $L^1$ at any positive remaining horizon.

The same argument has a converse interpretation. If the relevant code-rooted functionals are integrable, then the even terminal derivatives satisfy a factorial bound. Applying the estimate to both $f^*$ and $(\partial_{z_j}f)^*$ gives
\begin{equation*}
  \abs{\partial_{z_j}^r f(J_n\phi(y))}
  \leq
  C_yA_y^r\Gamma\pa{\frac r2+1}
\end{equation*}
for almost every $y$. Thus the formal directional Taylor series is entire with at most Gaussian growth. This Gevrey-$1/2$ necessity is stated in~\cref{thm:gevrey-necessity}. It shows how restrictive absolute integrability can be for a representation which repeatedly differentiates the nonlinearity along the tree.

These results do not say that the underlying PDE lacks a probabilistic representation. They say that one particular exact architecture can have an infinite absolute moment. The distinction can be seen on a single benchmark. For
\begin{equation}
  \partial_tu+\frac12u_{xx}
  +\eta\pa{e^{(u_x)^4}-1}=0,
  \qquad
  u(T,x)=a\cos x,
  \label{eq:intro-benchmark}
\end{equation}
the Nguwi--Penent--Privault functional is non-$L^1$ at every positive horizon, while the marked branching representation of Henry-Labord\`ere, Oudjane, Tan, Touzi, and Warin~\cite{HenryLabordereEtAl2019} has a finite second moment on an explicit positive interval. The precise statement is~\cref{thm:representation-dichotomy}. We use this example only to establish the first step of the paper's logic: formal exactness does not imply $L^1$, and the moment question cannot be read off from the PDE alone.

\subsection{The quadratic Hessian equation and patch regrouping}

The main part of the paper studies
\begin{equation}
  \partial_tv
  =
  \frac12v_{xx}+\lambda(v_{xx})^2
  \qquad\text{on }\Torus=\R/(2\pi\mathbb Z).
  \label{eq:intro-hessian-pde}
\end{equation}
Writing $z=v_{xx}$ gives the mild equation
\begin{equation}
  z(t)
  =
  P_t\phi''
  +\lambda\int_0^t
  \partial_x^2P_{t-s}\br{z(s)^2}\,ds.
  \label{eq:intro-hessian-mild}
\end{equation}
Repeated substitution produces planar binary Duhamel trees. Each internal vertex carries a Hessian transfer. We group maximal chains obtained by repeatedly following the left child into \emph{patches}. The finite-depth statement is elementary but important: replacing every such chain by its complete iterated Duhamel integral is an exact reindexing of the finite Picard expansion. No estimate and no infinite sum enter the argument. This is~\cref{thm:finite-patch-regrouping}.

The regrouping changes where one is allowed to take absolute values. Along a patch, several centered Hessian transfers can be composed first. In a bare Gaussian chain, conditioning on the total displacement replaces the product of second Hermite scores by one higher Hermite score. With spatial multipliers between the derivative transfers, the same cancellation is expressed by commutator expansions. The finite patch theorem itself is only a signed identity; the point is to make these cancellations available before absolute moments are estimated.

A deterministic semi-implicit iteration gives a small uniformly parabolic solution of~\eqref{eq:intro-hessian-pde}. We include it because it identifies the deterministic solution represented later and gives a direct self-consistent diffusion formula. It is not the main result. Under a uniform Schauder smallness condition, the iteration
\begin{equation*}
  \partial_tv_{n+1}
  =
  \pa{\frac12+\lambda z_n}v_{n+1,xx},
  \qquad
  z_{n+1}=v_{n+1,xx},
\end{equation*}
remains in a fixed ellipticity window and contracts in $H^{-1}$. The resulting theorem is stated in~\cref{thm:deterministic-iteration}. Its role is auxiliary: it supplies a deterministic solution theory on a regime which overlaps the probabilistic theorem below.

\subsection{Averaging before absolute values}

Let
\begin{equation*}
  X_{\alpha,T}=C^{\alpha/2,\alpha}([0,T]\times\Torus),
  \qquad
  (\cD f)(t)=\int_0^t\partial_x^2P_{t-s}f(s)\,ds.
\end{equation*}
The Hessian Duhamel operator is bounded on this parabolic H\"older scale in the form
\begin{equation*}
  \norm{\cD(fg)}_X
  \leq
  C_{\cD}(\alpha,T)\norm f_X\norm g_X.
\end{equation*}
For a finite rooted planar binary tree $\tau$, define its \emph{interior-averaged profile} recursively by
\begin{equation*}
  F_\bullet=P_\cdot\phi'',
  \qquad
  F_{[\tau_1,\tau_2]}
  =
  \lambda\cD(F_{\tau_1}F_{\tau_2}).
\end{equation*}
All continuous branch times, Brownian increments, and Hermite marks inside the decorated skeleton have been integrated out in $F_\tau$. If
\begin{equation}
  4\abs\lambda C_{\cD}(\alpha,T)
  \norm{P_\cdot\phi''}_{X_{\alpha,T}}<1,
  \label{eq:intro-cprime-smallness}
\end{equation}
then the Catalan count of binary trees and the bilinear bound above give
\begin{equation*}
  \sum_{\tau}\norm{F_\tau}_X<\infty.
\end{equation*}
Consequently the skeleton series is absolutely convergent, solves~\eqref{eq:intro-hessian-mild}, and may itself be randomized. If $S$ is sampled from any full-support probability mass function $\pi$ on the finite decorated skeletons, then
\begin{equation*}
  \widehat z(t,x)=\frac{F_S(t,x)}{\pi(S)}
\end{equation*}
belongs to $L^1$ and has expectation $z(t,x)$. This is Theorem C-prime, stated as~\cref{thm:cprime}.

The construction is intentionally ordered as follows: first average every continuous variable inside the skeleton, then take absolute values, then randomize the countable skeleton. One cannot instead begin from an unresolved infinite raw functional $H$ and write $\E[H\mid S]$, because ordinary conditional expectation already requires $H\in L^1$. The deterministic Duhamel recursion defines $F_S$ without this circularity.

\subsection{Raw-faithfulness is impossible for one fixed smooth datum}

The next result asks whether the same integrability can be obtained while retaining the canonical raw signed marked integrand. The class must be stated carefully. For a fixed finite raw marked genealogy, let $\mu$ denote its intrinsic finite signed measure after the positive proposal factors and their reciprocal compensators have been cancelled. Let $Q$ be any positive proposal dominating $\mu$. An estimator $Y$ is \emph{raw-faithful} if, after the canonical raw marks are exposed,
\begin{equation}
  \E_Q[Y\mid\text{raw marks}]
  =
  \frac{d\mu}{dQ}.
  \label{eq:intro-raw-faithful}
\end{equation}
The proposal may change the lifetime law, the genealogy distribution, the Gaussian law, or dependencies among these variables. Additional auxiliary randomness is also allowed. Condition~\eqref{eq:intro-raw-faithful} requires only that these changes do not move signed mass between distinct canonical raw marked states.

For one fixed smooth datum, this class has no $L^1$ estimator. The datum is built directly into the terminal Hessian profile,
\begin{equation}
  g_\varepsilon(x)
  =
  \varepsilon\br{
    \cos x+
    \sum_{m\geq m_0}
    \frac{\cos(N_mx)}{\sqrt{m!}}
  },
  \qquad
  N_m=K^m,
  \label{eq:intro-lacunary-datum}
\end{equation}
with $K\geq3$. For each fixed derivative order $k$,
\begin{equation*}
  \sum_m\frac{N_m^k}{\sqrt{m!}}<\infty,
\end{equation*}
so~\eqref{eq:intro-lacunary-datum} is $C^\infty$. It is the second derivative of a smooth periodic terminal datum $\phi_\varepsilon$.

The lower bound comes from a disjoint family of length-$m$ comb genealogies. Side leaves are projected onto frequency $1$ and the terminal distinguished leaf onto frequency $N_m$. Restrict every Hessian duration to
\begin{equation*}
  N_m^{-2}\leq r_j\leq\frac{h}{4m}.
\end{equation*}
This rectangle lies inside the chronological simplex. After absolute values, each retained centered Hessian mark contributes a factor of order
\begin{equation*}
  \int_{N_m^{-2}}^{h/(4m)}\frac{dr}{r}
  \asymp m.
\end{equation*}
The intrinsic total variation $\norm{\mu_m}_{\TV}$ of the length-$m$ comb therefore satisfies
\begin{equation}
  \norm{\mu_m}_{\TV}
  \gtrsim
  \varepsilon(C\abs\lambda\varepsilon)^m
  \frac{m^m}{\sqrt{m!}}.
  \label{eq:intro-raw-tv}
\end{equation}
The right-hand side is not summable. Since the comb cylinders are disjoint, conditional Jensen applied to~\eqref{eq:intro-raw-faithful} gives
\begin{equation*}
  \E_Q\abs Y
  \geq
  \sum_m\norm{\mu_m}_{\TV}
  =\infty.
\end{equation*}
This is~\cref{thm:raw-faithful-obstruction}. It is proposal invariant because the first moment of the Radon--Nikodym derivative is the total variation of the same intrinsic signed measure.

The amplitude $\varepsilon$ may be chosen arbitrarily small. In particular, the same fixed datum may be placed inside the C-prime regime~\eqref{eq:intro-cprime-smallness}. This gives the clean contrast of the paper:
\begin{equation*}
  \boxed{
  \begin{minipage}{0.82\textwidth}
  For the same small smooth terminal datum, complete interior averaging gives an unbiased $L^1$ representation, while every raw-faithful representation has infinite first moment.
  \end{minipage}}
\end{equation*}

\subsection{What the negative theorem does not say}

The raw-faithful theorem is not an impossibility theorem for randomness. This point belongs in the main statement of the programme because otherwise the scope is easy to misread.

There is first a vacuous counterexample. Start from the C-prime estimator and sample, in addition, an independent family of Gaussian variables labelled by the interior edges of the sampled skeleton. If the returned scalar is still $F_S/\pi(S)$, the simulation contains continuous random marks but their presence has no mathematical effect. Thus a condition which merely requires random variables to be present cannot distinguish raw retention from decorative randomness.

There are also nontrivial escape mechanisms. At one Hessian edge, the antithetic estimator
\begin{equation}
  \widetilde K_rF(x,Z)
  =
  \frac{\He_2(Z)}{2r}
  \br{F(x+\sqrt rZ)+F(x-\sqrt rZ)-2F(x)}
  \label{eq:intro-antithetic}
\end{equation}
still uses the sampled Gaussian mark and has expectation $\partial_x^2P_rF$, but conditional on $Z$ it is not the canonical one-sided raw integrand~\eqref{eq:intro-centered-hessian}. It therefore falls outside raw-faithfulness. Antithetic and ghost constructions of this kind are standard variance-reduction devices in branching methods; see, for example,~\cite{Warin2017}.

The surviving problem is therefore not whether random marks may appear. It is whether one can redistribute signed mass between raw marked states, while retaining nontrivial continuous randomness, enough to make the total variation summable.

\subsection{Coarsening signed patch measures}

The preceding distinction has a useful measure-theoretic formulation. For a finite decorated skeleton $\tau$, let $\mu_\tau$ be its intrinsic signed measure on the canonical raw interior-mark space $\cM_\tau$. Its total mass is the deterministic skeleton contribution $F_\tau(t,x)$. Let
\begin{equation*}
  \cC_\tau:\cM_\tau\longrightarrow\overline{\cM}_\tau
\end{equation*}
be a measurable coarsening and define
\begin{equation}
  \overline\mu_\tau
  =
  (\cC_\tau)_{\#}\mu_\tau.
  \label{eq:intro-coarsening}
\end{equation}
A canonical importance sampler for the coarsened state has absolute first moment
\begin{equation*}
  \sum_\tau\norm{\overline\mu_\tau}_{\TV}.
\end{equation*}
Thus the open problem becomes the construction of a nonconstant family of coarsenings for which this sum is finite.

Three points locate the problem. If $\cC_\tau$ is the identity,~\cref{thm:raw-faithful-obstruction} gives divergence for the datum~\eqref{eq:intro-lacunary-datum}. If $\cC_\tau$ is constant, then
\begin{equation*}
  \norm{(\cC_\tau)_{\#}\mu_\tau}_{\TV}
  =
  \abs{F_\tau(t,x)},
\end{equation*}
and Theorem C-prime proves summability under~\eqref{eq:intro-cprime-smallness}. More generally,
\begin{equation}
  \norm{(\cC_\tau)_{\#}\mu_\tau}_{\TV}
  \geq
  \abs{\mu_\tau(\cM_\tau)},
  \label{eq:intro-tv-optimality}
\end{equation}
so the constant map is total-variation optimal within the coarsening class. This optimality does not imply that it is the only finite choice.

A natural intermediate map keeps the total Gaussian displacement along a Hessian chain and integrates out only the Gaussian bridge coordinates. For a bare chain, this replaces a product of $m$ second Hermite scores by one $2m$-th Hermite score and changes factorial retained-mark growth into geometric growth after the time simplex is integrated. Whether an analogous nonconstant coarsening yields
\begin{equation*}
  \sum_\tau\norm{\overline\mu_\tau}_{\TV}<\infty
\end{equation*}
for the full spatially varying patch expansion is left open.

\subsection{Organization}

\Cref{sec:preliminaries} fixes the semigroup, signed-measure, tree, and coarsening notation. \Cref{sec:coding-tree} states the repeated-Hessian obstruction and Gevrey-$1/2$ necessity. \Cref{sec:dichotomy} gives the Nguwi--Penent--Privault versus HLOTW benchmark. \Cref{sec:patches} introduces the finite patch regrouping and records the auxiliary deterministic iteration. \Cref{sec:cprime} states the skeleton-averaged $L^1$ theorem. \Cref{sec:raw-faithful} states the fixed-datum raw-faithful obstruction and isolates the role of decorative and antithetic randomness. \Cref{sec:coarsening} formulates the remaining problem through pushforwards of the intrinsic signed patch measures.
